% !TeX program = xelatex
\documentclass[12pt,a4paper]{article}
\usepackage{fontspec}
\setmainfont{Times New Roman}
\usepackage{amsmath,amssymb}
\usepackage[margin=2cm,includeheadfoot]{geometry}
\usepackage{enumitem}
\usepackage{microtype}
\setlength{\parindent}{0pt}
\setlength{\parskip}{0.35em}
\setlist[enumerate]{leftmargin=*,itemsep=0.55em,topsep=0.3em}
\newcommand{\diem}[1]{\textit{(#1 điểm)}}

\begin{document}
\begin{center}
  {\large\bfseries SỞ GIÁO DỤC VÀ ĐÀO TẠO HÀ NỘI}\par
  {\Large\bfseries KỲ THI TUYỂN SINH VÀO LỚP 10 THPT}\par
  {\large\bfseries NĂM HỌC 2026--2027}\par
  \vspace{0.35em}
  {\large\bfseries ĐÁP ÁN VÀ LỜI GIẢI CHI TIẾT}\par
  Môn thi: Toán \qquad Ngày thi: 31/5/2026
\end{center}
\textit{Đáp án cho phần trích Đại số và toán thực tế trong đề; tổng điểm các ý được giữ lại theo thang điểm gốc là 7,0 điểm.}

\textbf{Câu I} \diem{1,5}
\begin{enumerate}
  \item Nhóm $[150;155)$ có tần số là số học sinh thuộc nhóm đó, do đó
  \[
    n_{[150;155)}=14.
  \]
  Tần số tương đối của nhóm là
  \[
    \frac{14}{50}=0{,}28=28\%.
  \]
  \textit{Kết luận: tần số bằng $14$, tần số tương đối bằng $28\%$.}

  \item Không gian mẫu gồm 6 kết quả đồng khả năng là lấy được một trong các quả bóng ghi số $1,2,3,4,5,6$. Biến cố $A$ gồm 3 kết quả thuận lợi: lấy được bóng ghi số $2,4$ hoặc $6$. Vậy
  \[
    P(A)=\frac{3}{6}=\frac12.
  \]
\end{enumerate}

\textbf{Câu II} \diem{1,5}\\
Đặt điều kiện như đề bài: $x>0$, $x\ne9$.
\begin{enumerate}
  \item Thay $x=25$ vào $A$:
  \[
    A=\frac{\sqrt{25}-4}{\sqrt{25}}=\frac{5-4}{5}=\frac15.
  \]

  \item Đặt $t=\sqrt{x}$. Khi đó $t>0$, $t\ne3$, và $x-9=(t-3)(t+3)$. Quy đồng hai phân thức:
  \begin{align*}
    B&=\frac{4(t+3)+t^2-7t-12}{(t-3)(t+3)}\\
     &=\frac{t^2-3t}{(t-3)(t+3)}
      =\frac{t}{t+3}
      =\frac{\sqrt{x}}{\sqrt{x}+3}.
  \end{align*}

  \item Theo biểu thức gốc $B$, kết quả phần (2) suy ra
  \[
    P=A\cdot B=\frac{\sqrt{x}-4}{\sqrt{x}}\cdot\frac{\sqrt{x}}{\sqrt{x}+3}
      =\frac{\sqrt{x}-4}{\sqrt{x}+3}.
  \]
  Đặt $t=\sqrt{x}>0$, $t\ne3$. Khi đó $P=(t-4)/(t+3)=1-7/(t+3)$. Vì $t+3>3$, ta có $-\frac43<P<1$. Nếu $P$ nguyên thì $P\in\{-1,0\}$.
  \begin{itemize}
    \item $P=0\iff t=4\iff x=16$.
    \item $P=-1\iff t+3=\frac72\iff t=\frac12\iff x=\frac14$.
  \end{itemize}
  Cả hai giá trị đều thỏa điều kiện. Vậy $x=16$ hoặc $x=\frac14$.
\end{enumerate}

\textbf{Câu III} \diem{2,5}
\begin{enumerate}
  \item Gọi $x$ (chiếc áo) là số áo theo kế hoạch tổ may trong mỗi ngày, với $x>0$. Trong 3 ngày đầu tổ may được $3x$ chiếc; trong 7 ngày tiếp theo tổ may được $7(x+5)$ chiếc. Theo đề bài,
  \[
    3x+7(x+5)=335
    \iff 10x+35=335
    \iff x=30.
  \]
  Vậy theo kế hoạch, mỗi ngày tổ may 30 chiếc áo.

  \item Gọi $r$ và $c$ lần lượt là số bông hoa hồng và hoa cúc. Ta có hệ
  \[
    \begin{cases}
      r+c=25,\\
      8r+6c=180.
    \end{cases}
  \]
  Từ $c=25-r$, thay vào phương trình tiền:
  \[
    8r+6(25-r)=180
    \iff 2r=30
    \iff r=15.
  \]
  Suy ra $c=25-15=10$. Vậy người đó mua 15 bông hoa hồng và 10 bông hoa cúc.

  \item Theo hệ thức Viète cho $x^2-3x+1=0$,
  \[
    x_1+x_2=3,\qquad x_1x_2=1.
  \]
  Do $x_1x_2=1$, ta có $\dfrac1{x_1}=x_2$ và $\dfrac1{x_2}=x_1$. Vì thế
  \begin{align*}
    Q&=\frac{3x_2-1}{x_1}+\frac{3x_1}{x_2}-x_1\\
     &=3x_2^2-x_2+3x_1^2-x_1\\
     &=3\bigl((x_1+x_2)^2-2x_1x_2\bigr)-(x_1+x_2)\\
     &=3(9-2)-3=18.
  \end{align*}
\end{enumerate}

\textbf{Câu IV} \textit{(ý 1 của đề gốc; phần giữ lại: 1,0 điểm)}
\begin{enumerate}
  \item Xô có bán kính đáy $r=12\,\mathrm{cm}$, chiều cao $h=25\,\mathrm{cm}$.
  \begin{enumerate}[label=\alph*)]
    \item Diện tích xung quanh hình trụ là
    \[
      S_{\mathrm{xq}}=2\pi rh\approx2\cdot3{,}14\cdot12\cdot25=1884\,\mathrm{cm}^2.
    \]
    \item Thể tích của xô là
    \[
      V=\pi r^2h\approx3{,}14\cdot12^2\cdot25=11304\,\mathrm{cm}^3=11{,}304\,\mathrm{lít}.
    \]
    Mỗi lần chỉ múc lượng nước bằng $80\%$ thể tích xô, tức là
    \[
      0{,}8\cdot11{,}304=9{,}0432\,\mathrm{lít}.
    \]
    Sau 16 lần, lượng nước múc được là $16\cdot9{,}0432=144{,}6912$ lít, chưa đủ 150 lít; sau 17 lần là $17\cdot9{,}0432=153{,}7344$ lít, đủ đổ đầy bể. Vậy cần múc ít nhất 17 xô.
  \end{enumerate}
\end{enumerate}

\textbf{Câu V} \diem{0,5}\\
Gọi $n$ là số công nhân và $d$ là số ngày thuê xưởng. Theo cách lập kế hoạch hoàn thành đúng 1000 sản phẩm,
\[
  5nd=1000\quad\Longleftrightarrow\quad nd=200,
\]
trong đó $n,d$ là các số nguyên dương. Tính theo triệu đồng, tổng chi phí là
\[
  C=3d+n=\frac{600}{n}+n.
\]
Các ước dương của 200 và chi phí tương ứng là
\[
\begin{array}{c|rrrrrrrrrrrr}
 n&1&2&4&5&8&10&20&25&40&50&100&200\\ \hline
 C&601&302&154&125&83&70&50&49&55&62&106&203
\end{array}
\]
Do đó chi phí nhỏ nhất là 49 triệu đồng, đạt được khi $n=25$ và $d=8$. Vậy công ty nên điều động 25 công nhân và thuê kho xưởng trong 8 ngày.

\begin{center}\textbf{HẾT}\end{center}
\end{document}
